\textbf{Phần 1. Trắc nghiệm lựa chọn}
\setcounterpageref{ex}{0}

\begin{ex} %[5H8N47]
$1\text{ m}^3$ bằng bao nhiêu đề-xi-mét khối?
\choice
{$10\text{ dm}^3$}
{\True $1\,000\text{ dm}^3$}
{$100\text{ dm}^3$}
{$1\,000\,000\text{ dm}^3$}
\loigiai{
$1\text{ m}^3 = 1\,000\text{ dm}^3$.\\
\textbf{Đáp án của câu hỏi là B}
}
\end{ex}

\begin{ex} %[5H8N47]
$1\text{ m}^3$ bằng bao nhiêu xăng-ti-mét khối?
\choice
{$1\,000\text{ cm}^3$}
{$10\,000\text{ cm}^3$}
{\True $1\,000\,000\text{ cm}^3$}
{$100\,000\text{ cm}^3$}
\loigiai{
$1\text{ m}^3 = 1\,000\,000\text{ cm}^3$.\\
\textbf{Đáp án của câu hỏi là C}
}
\end{ex}

\begin{ex} %[5H8N47]
Số “Hai phẩy không năm mét khối” được viết là
\choice
{$205\text{ m}^3$}
{$2{,}5\text{ m}^3$}
{$2{,}005\text{ m}^3$}
{\True $2{,}05\text{ m}^3$}
\loigiai{
Viết là $2{,}05\text{ m}^3$.\\
\textbf{Đáp án của câu hỏi là D}
}
\end{ex}

\begin{ex} %[5H8H47]
Đổi $4{,}8\text{ m}^3$ ra đề-xi-mét khối ta được
\choice
{$48\text{ dm}^3$}
{$480\text{ dm}^3$}
{\True $4\,800\text{ dm}^3$}
{$48\,000\text{ dm}^3$}
\loigiai{
$4{,}8 \times 1\,000 = 4\,800\text{ dm}^3$.\\
\textbf{Đáp án của câu hỏi là C}
}
\end{ex}

\begin{ex} %[5H8H47]
Đổi $2\,500\text{ dm}^3$ ra mét khối ta được
\choice
{$0{,}25\text{ m}^3$}
{$25\text{ m}^3$}
{$250\text{ m}^3$}
{\True $2{,}5\text{ m}^3$}
\loigiai{
$2\,500 : 1\,000 = 2{,}5\text{ m}^3$.\\
\textbf{Đáp án của câu hỏi là D}
}
\end{ex}

\begin{ex} %[5H8H47]
$1\text{ m}^3$ nước tương ứng với bao nhiêu lít nước?
\choice
{$100\text{ l}$}
{\True $1\,000\text{ l}$}
{$10\,000\text{ l}$}
{$1\text{ l}$}
\loigiai{
Vì $1\text{ m}^3 = 1\,000\text{ dm}^3$ và $1\text{ dm}^3 = 1\text{ l}$ nên $1\text{ m}^3 = 1\,000\text{ l}$.\\
\textbf{Đáp án của câu hỏi là B}
}
\end{ex}

\begin{ex} %[5H8V47]
Một bể nước hình hộp chữ nhật chứa được tối đa $3{,}6\text{ m}^3$ nước. Dung tích của bể nước đó là bao nhiêu lít?
\choice
{\True $3\,600\text{ l}$}
{$360\text{ l}$}
{$36\,000\text{ l}$}
{$36\text{ l}$}
\loigiai{
$3{,}6\text{ m}^3 = 3{,}6 \times 1\,000 = 3\,600\text{ l}$.\\
\textbf{Đáp án của câu hỏi là A}
}
\end{ex}

\begin{ex} %[5H8V47]
Tính giá trị biểu thức: $4\text{ m}^3 - 1\,200\text{ dm}^3 \times 2$.
\choice
{$2\text{ m}^3$}
{$2{,}8\text{ m}^3$}
{\True $1{,}6\text{ m}^3$}
{$1{,}8\text{ m}^3$}
\loigiai{
$1\,200\text{ dm}^3 \times 2 = 2\,400\text{ dm}^3 = 2{,}4\text{ m}^3$.\\
$4 - 2{,}4 = 1{,}6\text{ m}^3$.\\
\textbf{Đáp án của câu hỏi là C}
}
\end{ex}

\begin{ex} %[5H8V47]
Một xe téc chở $15\text{ m}^3$ xăng. Người ta đã bơm từ xe téc vào các bồn chứa $8\,500\text{ l}$ xăng. Số mét khối xăng còn lại trên xe là
\choice
{\True $6{,}5\text{ m}^3$}
{$7\text{ m}^3$}
{$6\text{ m}^3$}
{$5{,}5\text{ m}^3$}
\loigiai{
Đổi $8\,500\text{ l} = 8{,}5\text{ m}^3$.\\
Số xăng còn lại: $15 - 8{,}5 = 6{,}5\text{ (m}^3\text{)}$.\\
\textbf{Đáp án của câu hỏi là A}
}
\end{ex}

\begin{ex} %[5H8V47]
Điền dấu thích hợp: $2\text{ m}^3\; 85\text{ dm}^3 \;\dots\; 2{,}85\text{ m}^3$.
\choice
{$>$}
{\True $<$}
{$=$}
{Không so sánh được}
\loigiai{
$2\text{ m}^3\; 85\text{ dm}^3 = 2{,}085\text{ m}^3 < 2{,}85\text{ m}^3$.\\
\textbf{Đáp án của câu hỏi là B}
}
\end{ex}

\begin{ex} %[5H8V47]
Một căn phòng có thể tích không khí là $72\text{ m}^3$. Biết chiều cao phòng là $3\text{ m}$. Diện tích nền phòng học là
\choice
{\True $24\text{ m}^2$}
{$216\text{ m}^2$}
{$36\text{ m}^2$}
{$48\text{ m}^2$}
\loigiai{
Diện tích nền phòng: $72 : 3 = 24\text{ (m}^2\text{)}$.\\
\textbf{Đáp án của câu hỏi là A}
}
\end{ex}

\begin{ex} %[5H8C47]
Một bể bơi có thể tích $180\text{ m}^3$. Người ta dùng máy bơm có công suất $15\text{ m}^3/\text{giờ}$ để bơm nước vào bể. Biết trong bể đã có sẵn $30\text{ m}^3$ nước. Thời gian để bơm đầy bể là
\choice
{$12\text{ giờ}$}
{$8\text{ giờ}$}
{$14\text{ giờ}$}
{\True $10\text{ giờ}$}
\loigiai{
Lượng nước cần bơm thêm: $180 - 30 = 150\text{ (m}^3\text{)}$.\\
Thời gian bơm đầy bể: $150 : 15 = 10\text{ (giờ)}$.\\
\textbf{Đáp án của câu hỏi là D}
}
\end{ex}

\textbf{Phần 2. Trắc nghiệm đúng sai}
\setcounterpageref{ex}{0}

\begin{ex} %[5H8H47]
Xét tính đúng sai của các khẳng định sau:
\choiceTF
{\True $1\text{ m}^3 = 1\,000\text{ dm}^3 = 1\,000\,000\text{ cm}^3$.}
{$0{,}75\text{ m}^3 = 75\text{ dm}^3$.}
{\True $2\text{ m}^3\; 50\text{ dm}^3 = 2{,}05\text{ m}^3$.}
{\True $12\,500\text{ dm}^3 = 12{,}5\text{ m}^3$.}
\loigiai{
\begin{itemize}
    \item $1\text{ m}^3 = 1000\text{ dm}^3 = 1000000\text{ cm}^3$. Đúng.
    \item $0{,}75 \times 1000 = 750\text{ dm}^3 \neq 75\text{ dm}^3$. Sai.
    \item $2 + 50 : 1000 = 2{,}05\text{ m}^3$. Đúng.
    \item $12500 : 1000 = 12{,}5\text{ m}^3$. Đúng.
\end{itemize}
\begin{itemize}
    \item a) \textbf{Đúng}.
    \item b) \textbf{Sai}.
    \item c) \textbf{Đúng}.
    \item d) \textbf{Đúng}.
\end{itemize}
}
\end{ex}

\begin{ex} %[5H8V47]
Một bể nước dạng hình hộp chữ nhật có thể tích $12\text{ m}^3$. Hiện lượng nước trong bể chiếm $65\%$ thể tích bể. Xét tính đúng sai:
\choiceTF
{\True Dung tích toàn bộ của bể là $12\,000\text{ l}$.}
{\True Lượng nước hiện có trong bể là $7{,}8\text{ m}^3$.}
{Lượng nước cần bơm thêm để đầy bể là $5\,000\text{ l}$.}
{\True Lượng nước cần bơm thêm để đầy bể là $4\,200\text{ l}$.}
\loigiai{
\begin{itemize}
    \item Dung tích: $12 \times 1000 = 12000\text{ l}$. Đúng.
    \item Nước hiện có: $12 \times 65\% = 7{,}8\text{ m}^3 = 7800\text{ l}$. Đúng.
    \item Cần thêm: $12000 - 7800 = 4200\text{ l} \neq 5000\text{ l}$. Ý 3 Sai, Ý 4 Đúng.
\end{itemize}
\begin{itemize}
    \item a) \textbf{Đúng}.
    \item b) \textbf{Đúng}.
    \item c) \textbf{Sai}.
    \item d) \textbf{Đúng}.
\end{itemize}
}
\end{ex}

\textbf{Phần 3. Trả lời ngắn}
\setcounterpageref{ex}{0}

\begin{ex} %[5H8H47]
Đổi $5{,}25\text{ m}^3$ ra đề-xi-mét khối.
\shortans{5250}
\loigiai{
\textbf{Đáp số:} $5250$\\
$5{,}25 \times 1\,000 = 5\,250\text{ dm}^3$.
}
\end{ex}

\begin{ex} %[5H8H47]
Đổi $3\,400\text{ dm}^3$ ra mét khối (dùng dấu phẩy nếu có).
\shortans{3,4}
\loigiai{
\textbf{Đáp số:} $3{,}4$\\
$3\,400 : 1\,000 = 3{,}4\text{ m}^3$.
}
\end{ex}

\begin{ex} %[5H8V47]
Một bể bơi có thể tích chứa là $45\text{ m}^3$. Hỏi bể bơi đó chứa được bao nhiêu nghìn lít nước?
\shortans{45}
\loigiai{
\textbf{Đáp số:} $45$\\
$45\text{ m}^3 = 45\,000\text{ l} = 45$ nghìn lít.
}
\end{ex}

\begin{ex} %[5H8V47]
Thực hiện phép tính: $8{,}5\text{ m}^3 - 3\,700\text{ dm}^3$. Kết quả bằng bao nhiêu mét khối?
\shortans{4,8}
\loigiai{
\textbf{Đáp số:} $4{,}8$\\
$8{,}5 - 3{,}7 = 4{,}8\text{ m}^3$.
}
\end{ex}

\begin{ex} %[5H8V47]
Một lớp học hình hộp chữ nhật có chiều dài $8\text{ m}$, chiều rộng $6\text{ m}$, chiều cao $3{,}5\text{ m}$. Tính thể tích không khí trong phòng học theo mét khối.
\shortans{168}
\loigiai{
\textbf{Đáp số:} $168$\\
$8 \times 6 \times 3{,}5 = 168\text{ m}^3$.
}
\end{ex}

\begin{ex} %[5H8C47]
Một bể nước hình hộp chữ nhật chứa $18\text{ m}^3$ nước. Mỗi ngày gia đình dùng hết $750\text{ l}$ nước. Hỏi sau bao nhiêu ngày thì gia đình dùng hết số nước trong bể?
\shortans{24}
\loigiai{
\textbf{Đáp số:} $24$\\
$18\text{ m}^3 = 18\,000\text{ l} \Rightarrow 18\,000 : 750 = 24\text{ ngày}$.
}
\end{ex}
